\documentclass[12pt]{report}

\textheight 22cm
\textwidth 15.5cm
\oddsidemargin 0pt\evensidemargin 0pt
%\oddsidemargin 14pt\evensidemargin 0pt
%\topmargin -40pt
\topmargin-30pt
%\bottommargin0pt
\def\baselinestretch{1.1}
%\addtolength{\parskip}{1ex}
\jot=.5ex
%\parskip = 0.02in


\setlength\arraycolsep{2pt}


\usepackage{amssymb}
\usepackage{amsmath,bm}
\usepackage{amssymb}
\usepackage{graphicx}
\usepackage{amsfonts}         
\usepackage{fancybox}   

%\usepackage[numbers]{natbib}

\usepackage{enumitem}

\usepackage{slashed}




\usepackage[usenames,dvipsnames]{xcolor}%http://en.wikibooks.org/wiki/LaTeX/Colors

\definecolor{darkgreen}{rgb}{0,0.4,0}
\definecolor{darkred}{rgb}{0.4,0,0}
\definecolor{darkblue}{rgb}{0,0,0.4}
\definecolor{lightblue}{rgb}{.6,.6,0.9}
\newcommand{\cobl}{\color{darkblue}}

\newcommand{\cor}{\color{red}}
\newcommand{\cog}{\color{darkgreen}}
\newcommand{\cob}{\color{black}}

\definecolor{uglybrown}{rgb}{0.8,  0.7,  0.5}


\def\nd{{ \vphantom{\dagger}}}

\def\dd{\text{d}}
\def\dbar{{\mathchar'26\mkern-12mu \dd}} 

\def\ii{{\bf i}}
\def\Ione{\mathbb{I}}
\def\UU{{\bf U}}
\def\HH{{\bf H}}
\def\pp{{\bf p}}
\def\aa{{\bf a}}
\def\qq{{\bf q}}
\def\eps{\epsilon}
\def\half{{1\over 2}}
\def\Tr{{{\rm Tr~ }}}
\def\tr{{\rm tr}}
\def\Re{{\rm Re\hskip0.1em}}
\def\Im{{\rm Im\hskip0.1em}}
\def\ppi{\boldsymbol{\pi}}
\def\pphi{\boldsymbol{\phi}}
%\def\pphi{\phi}
\def\grad{\vec \nabla}
\def\vB{\vec B}
\def\vE{\vec E}
\def\vA{\vec A}
\def\vAA{ \vec{\bf A}}
\def\vEE{{{\vec {\bf E}}}}
\def\vBB{{\vec {\bf B}}}
\def\CO{{\cal O}}
\def\CL{{\cal L}}
\def\IZ{\mathbb{Z}}
\def\CN{{\cal N}}
\def\CM{{\cal M}}
\def\IR{\mathbb{R}}
\def\CD{{\cal D}}
\def\gO{{\mathsf{O}}}


\def\ie{{\it i.e.}}
\def\eg{{\it e.g.}}

\def\bra#1{\left\langle#1\right|}
\def\ket#1{\left|#1\right\rangle}
\def\bbra#1{{\langle\langle}#1|}
\def\kket#1{|#1\rangle\rangle}
\def\vev#1{\left\langle{#1}\right\rangle}

\def\ketbra#1#2{ | #1 \rangle\hskip-2pt\langle #2|}



\def\be{\begin{equation}}
\def\ee{\end{equation}}
\def\({\left(}
\def\){\right)}


\usepackage{tikz}
\usetikzlibrary{decorations.markings}
\usetikzlibrary{arrows,decorations.pathmorphing,backgrounds,positioning,fit,petri,automata,shadows,calendar,mindmap, graphs}

% from flip tanedo.
% Based on:
% http://www.texample.net/tikz/examples/feynman-diagram/

\tikzset{
	% >=stealth', %% more traditional arrows, I don't like them
    vector/.style={decorate, decoration={snake}, draw},
    fermion/.style={postaction={decorate},
        decoration={markings,mark=at position .55 with {\arrow{>}}}},
    fermionbar/.style={draw, postaction={decorate},
        decoration={markings,mark=at position .55 with {\arrow{<}}}},
    fermionnoarrow/.style={},
    gluon/.style={decorate,
        decoration={coil,amplitude=4pt, segment length=5pt}},
    scalar/.style={dashed, postaction={decorate},
        decoration={markings,mark=at position .55 with {\arrow{>}}}},
    scalarbar/.style={dashed, postaction={decorate},
        decoration={markings,mark=at position .55 with {\arrow{<}}}},
    scalarnoarrow/.style={dashed,draw},
%
%%% 	Special vectors (when you need to fine-tune wiggles)
%	provector/.style={decorate, decoration={snake,amplitude=2.5pt}, draw},
%	antivector/.style={decorate, decoration={snake,amplitude=-2.5pt}, draw},
%	    electron/.style={draw=black, postaction={decorate},
%        decoration={markings,mark=at position .55 with {\arrow[draw=black]{>}}}},
%	bigvector/.style={decorate, decoration={snake,amplitude=4pt}, draw},
	vectorscalar/.style={loosely dotted,draw=black, postaction={decorate}},
}




\newcommand{\bea}{\begin{eqnarray}}
\newcommand{\eea}{\end{eqnarray}}


\def\parfig#1#2{
\parbox{#1\textwidth}
{\includegraphics[width=#1\textwidth]{#2}}
}



%%from  Rudro Rana Biswas
\usepackage[pagebackref,  % this puts links to the page numbers where refs appear
%pdftex, 
bookmarks={false}, pdfauthor={John McGreevy}, pdftitle={Yay, physics!}]{hyperref}
\hypersetup{colorlinks=true, 
linkcolor=BrickRed, 
citecolor=Violet, 
filecolor=OliveGreen, 
urlcolor=RoyalBlue, 
filebordercolor={.8 .8 1}, 
urlbordercolor={.8 .8 0}}%http://en.wikibooks.org/wiki/LaTeX/Hyperlinks


%\usepackage{mathtools} % for inclusion arrow \xhookrightarrow{}


\renewcommand{\theequation}{\arabic{equation}}
\newif{\ifeq}           % defines a new condition @eq tested by the conditional \ifeq
\eqtrue                 % if uncommented, declares @eq to be true
\eqfalse              % if uncommented, declares @eq to be false
%                                %
%                                % to use this, wrap text with the conditional, eg:
%                                %
%                                % \ifeq
%                                % SHOW THIS IFF \eqtrue HAS BEEN DECLARED
%                                % \fi
%
\def\answer#1
{
\ifeq
\textcolor{darkblue}{#1}
\fi
}


\begin{document}
\begin{center}

University of California at San Diego -- 
Department of Physics --
Prof.~John McGreevy

{\Large\bf  Physics 215B QFT Winter 2025}\\
{\Large\bf Assignment 6 \answer{--~~~ Solutions} }
\end{center}


\noindent
%Posted September 17, 2015 
\hfill {\bf Due  11:59pm Tuesday, February 18, 2025}




\bigskip
\hrule
\begin{enumerate}
\item {\bf Brain-warmer.}
Check that $ \(\Delta_T\)^\mu_\rho\equiv 
\delta^\mu_\rho - {q^\mu q_\rho\over q^2} $ 
is a projector onto momenta transverse to $q^\rho$.
This requires showing both that $ \Delta q = 0 $ and that $\Delta^2 = \Delta$.

\def\CA{{\cal A}}

\item {\bf A better formula for the superficial degree of divergence.}
[Thanks to Haoran Sun for suggesting this formula.]

Starting from the definition of an (amputated!) amplitude $\CA$ 
from a connected Feynman diagram, show that its superficial degree of divergence is 
\be 
D_\CA = 
D + \sum_{ \{g\}}  [g] V_g(\CA) - \sum_{ \{ f\} } E_f(\CA) [f] \ee
where $\{g\}$ is the set of coupling constants and $\{ f \}$ is the set of fields,
$V_g(\CA)$ is the number of vertices of the coupling $g$ in the diagram $\CA$,
and $E_f(\CA)$ is the number of external $f$ fields.  
For example, for the Yukawa theory you studied on a previous homework, this formula reduces to 
\be 
D_\CA = 
D +  [g] V_g + [y] V_y -B_E [\phi] - F_E [\phi] \ee
where $B_E\equiv E_\phi$ is the number of external scalars and $F_E \equiv E_\psi$ is the number of external fermions in the diagram.  
If you prefer, for definiteness, you could just show the formula for this case.  



\item {\bf Symmetry is attractive.}
 Consider a field theory in $D=3+1$ with two scalar fields
 with the same mass
which interact via the interaction 
$$ V = - { g \over 4!} \( \phi_1^4 + \phi_2^4 \) - { 2 \lambda \over 4!} \phi_1^2 \phi_2^2 .$$

\begin{enumerate}
\item Show that when $ \lambda= g$ the model 
possesses an $\gO(2)$ symmetry.


\item Will you need a counterterm of the form $ \phi_1 \phi_2$ 
or $ \phi_1 \Box \phi_2$ (for general $g, \lambda$)?  If not, why not?  



\item  Renormalize the theory to one loop order 
by regularizing (for example with a euclidean momentum cutoff or Pauli Villars),
adding the necessary counterterms, and imposing a renormalization
condition on the propagators  (consider the case where the scalars are both massless) and $ 2 \to 2 $ scattering amplitudes
at some values of the kinematical variables $s_0, t_0, u_0$.
Feel free to re-use our results from $\phi^4$ theory where appropriate.
%at some energy $\sqrt{s_0}$ and momentum transfer $\sqrt{t_0}$.



\item Consider the limit of low energies, \ie~when $ s_0, t_0, u_0 \ll \Lambda^2$ where $\Lambda$ is the cutoff scale.
Tune the location of the poles in both propagators to $p^2=0$.
Show that the coupling goes to the $\gO(2)$-symmetric value if it starts nearby (nearby means $ \lambda/g < 3$).  
(That is, show that at fixed physical coupling, 
the ratio of bare couplings $\lambda/g \to 1$ as we take the cutoff to infinity.)
A nice way to organize this is by computing the beta function for the coupling $\lambda/g$.


\end{enumerate}
\item {\bf Yes, please, gauge invariance.}  Verify for yourself that the one-loop vacuum polarization amplitude in QED (when computed using either the improved Pauli-Villars regulator or dim reg) satisfies the Ward identity, \ie~is proportional to 
$ q^\mu q^\nu - \eta^{\mu\nu} q^2 $.  It's up to you how much of this to hand in.


\item {\bf Soft gravitons?} [optional]
Photons are massless, and this means 
that the cross sections we measure 
actually include soft ones that we don't detect.
If we don't include them we get IR-divergent nonsense.

Gravitons are also massless.  Do we need to worry about them in the same way?
Here we'll sketch some hints for how to think about this question.  
%Warning: I just wrote these problems and haven't debugged them carefully yet.
%If something seems questionable to you, please ask me about it.


\begin{enumerate}
\item Consider the action
$$ S_0[h_{\mu\nu}] = \int d^4 x  \half h_{\mu\nu} \Box h^{\mu\nu} .$$
This is a kinetic term for (too many polarizations of a) two-index symmetric-tensor field $h_{\mu\nu} = h_{\nu\mu}$
(which we'll think of as a small fluctuation of the metric about flat space: $ g_{\mu\nu} \simeq \eta_{\mu\nu} + h_{\mu\nu}$, 
and this is where the coupling below comes from).
Like with the photon, we'll rely on the couplings to matter to keep
unphysical polarizations from being made.
Write the propagator for $h$.
We still raise and lower indices with $\eta_{\mu\nu}$.
\footnote{A warning:
%\answer{
I've done two misdeeds in the statement of this problem.
First, the Einstein-Hilbert term is 
$ \int d^4 x {1\over 8 \pi G_N}  \sqrt{g } R  = \int d^4 x {1\over 8 \pi G_N} ( \partial h)^2 + ...$
-- it has a factor of $G_N$ in front of it.  $R$ has units of $1 \over \text{length}^2$, 
and $g$ is dimensionless, so $G_N$ has units of $\text{length}^2$ -- it is 
$  8 \pi G_N = {1\over M_{\text{Pl}}^2} $, where $M_{\text{Pl}}$ is the Planck mass.
I've absorbed a factor of $ \sqrt{G_N}$ into $h$ 
so that the coefficient of the kinetic term is unity.
Second, the $(\partial h)^2$ here involves various index contractions, which I haven't written.
Some gauge fixing (de Donder gauge) is required to arrive at the simple expression I wrote above,
and one more thing -- the $h_{\mu\nu}$ I've written is actually the `trace-reversed' 
graviton field
$$ \bar h_{\mu\nu} \equiv h_{\mu\nu} - \half h \eta_{\mu\nu} $$
where $ h \equiv \eta^{\mu\nu} h_{\mu\nu}$ is the trace.
(I didn't write the bar.)
For the details of this, which are not needed for this problem, see chapter 10 of 
\htmladdnormallink{my GR notes}{http://mcgreevy.physics.ucsd.edu/f13/225A-lectures.pdf}.
}



\item Couple the graviton to the electron field via 
\bea\label{eq:couple-to-gravity} S_G &=& \int d^4 x~G h^{\mu\nu} T_{\mu\nu}  
\cr 
T_{\mu\nu} &\equiv& \bar \psi \( \gamma_\mu \partial_\nu  + \gamma_\nu \partial_\mu  \) \psi .\eea
What are the engineering dimensions of the coupling constant $G$? 
What is the new Feynman rule?
%(Note that this coupling is only the leading one in a more complete description.)


\item  Draw a (tree level) Feynman diagram that describes the creation
of gravitational radiation from an electron as 
a result of its acceleration from the absorption of a photon
($ e \gamma \to e h $).
Evaluate it if you dare.
Estimate or calculate the cross section (hint: use dimensional analysis).


\item Now the main event: 
study the one-loop diagram by which the 
graviton corrects the QED vertex.
Is it IR divergent?  If not, why not?



%\answer{
%The extra powers of the momentum in the numerator from the derivative coupling
%prevent an IR divergence.
%}

\item If you get stuck on the previous part,
replace the graviton field by a massless scalar $\pi(x)$.
Compare the case of derivative coupling $ \lambda \partial_\mu \pi \bar \psi \gamma^\mu \psi $
with the more direct Yukawa coupling $ y \pi \bar \psi\psi $.
[Warning: though this example has some similarities with the graviton case, the conclusion is different.]



\item Quite a bit about the form of the coupling of gravity to matter
is determined by the demand of coordinate invariance.
This plays a role like gauge invariance in QED.  
Acting on the small fluctuation, the transformation is
$$ h_{\mu\nu}(x)  \to h_{\mu\nu}(x) + \partial_\mu \lambda_\nu(x)  + \partial_\nu \lambda_{\mu}(x) .$$
What condition does the invariance under this (infinitesimal) transformation impose
on the object $T_{\mu\nu}$ appearing in \eqref{eq:couple-to-gravity}.


\end{enumerate}




\item {\bf Equivalent photon approximation.} [optional]
Consider a process in which very high-energy electrons scatter off a target.  
At leading order in $\alpha$, the electron line is connected
to the rest of the diagram by a single photon propagator.  
If the initial and final energies of the electron are $E$ and $E'$, the photon
will carry momentum $q$ with $q^2 = - 2 EE' (1-\cos\theta)$ 
(ignoring the electron mass $m \ll E$).  In the limit
of forward scattering ($\theta \to 0$), we have $q^2 \to 0$, 
so the photon approaches its mass shell.  In this problem, we ask: To what extent
can we treat it as a real photon?

\begin{enumerate}
\item The matrix element for the scattering process can be written as 
$$ \CM = - \ii e \bar u(p') \gamma^\mu u(p) { - \ii \eta_{\mu\nu} \over q^2 } \hat \CM^\nu(q) $$
where $\hat \CM^\nu$ represents the coupling of the virtual photon to the target.  
Let $ q = (q^0, \vec q)$ and define $ \tilde q = (q^0, - \vec q)$.  
The contribution to the amplitude from the electron line can be parametrized as 
$$ \bar u(p') \gamma^\mu u(p) = A q^\mu + B \tilde q^\mu + C \eps_1^\mu + D \eps_2^\mu $$
where $\eps_\alpha$ are unit vectors transverse to $\vec q$.   
Show that $B$ is at most of order $\theta^2$ (dot it with $q$), so we can ignore 
it at leading order in an expansion about forward scattering.
Why do we not care about the coefficient $A$?

%\answer{I will wait to post the solutions to this problem, since
%some folks said they would like more time to think about it.}


%
\item Working in the frame with $p = (E,0,0,E)$, compute
$$ \bar u(p') \gamma \cdot \eps_\alpha u(p) $$
explicitly using massless electrons, where $\bar u$ and $u$ are spinors of definite helicity, 
and $\eps_{\alpha = \parallel, \perp }$ are unit vectors parallel and perpendicular to the plane of 
scattering.  Keep only terms through order $\theta$.  
Note that for $\eps_\parallel$, the (small) $\hat 3$ component matters.



\item Now write the expression for the electron scattering cross section, in terms of 
$| \hat \CM^\mu|^2$ and the integral over phase space of the target.  
This expression must be integrated over the final electron momentum $\vec p'$.   The integral
over $p^{3'}$ is an integral over the energy loss of the electron.  Show that the integral over $p'_\perp$ 
diverges logarithmically as $ p'_\perp$ or $\theta \to 0$.


\item The divergence as $\theta \to 0$ 
is regulated by the electron mass (which we've ignored above).
Show that reintroducing the electron mass in the expression 
$$ q^2 = - 2 (EE' - p p' \cos \theta) + 2 m^2 $$
cuts off the divergence and gives a factor of $\log\(s/m^2\)$ in its place.



\item Assembling all the factors, and assuming that the target cross sections
are independent of photon polarization, show that the largest part of the electron-target
cross section is given by considering the electron to be the source of a beam of real photons
with energy distribution given by
$$ N_\gamma(x) dx = { dx \over x} {\alpha \over 2\pi} ( 1 + (1-x)^2 ) \log {s \over m^2 } $$
where $ x \equiv E_\gamma/E$.  
This is the Weisz\"acker-Williams equivalent photon approximation.  It is a precursor 
to the theory of jets and partons in QCD. 


\end{enumerate}





\end{enumerate}
\end{document}
